\scanpage{221}
\deleted{\mspage{110}}

\unsure{if|where} the corresponding \unsure{outer} operations of $\Gamma_{g,\nu}$ on $\Pi_{g,\nu}$ are those deduced by passage to the quotient from those of \deleted{of the \unsure{inclusion} $\Gamma_{g,\nu}\subset$}
$A_{g,\nu}\simeq\operatorname{Aut}(U_{g,\nu})$ on $\pi_1(U_{g,\nu},a_{g,\nu})=\Pi_{g,\nu}$.

The extension structure (41) is deduced from \unsure{(41)|(43)} by inverse image under the inclusion
\[
\Gamma_{g,\nu}^{!}\longrightarrow\Gamma_{g,\nu}.
\]

\uline{Suppose} we are in the anabelian case $2g+\nu\geq3$.
As for every extension structure by a kernel with trivial center, there is a canonical homomorphism into the canonical extension structure associated to $\Pi_{g,\nu}$
\[
(44)\quad
\diagram{\begin{tikzcd}[ampersand replacement=\&,column sep=small]
1 \arrow[r] \& \Pi_{g,\nu} \arrow[r,hook] \arrow[d,equal] \& \Gamma'_{g,\nu+1} \arrow[r] \arrow[d] \& \Gamma_{g,\nu} \arrow[r] \arrow[d] \& 1 \\
1 \arrow[r] \& \Pi_{g,\nu} \arrow[r] \& \operatorname{Aut}(\Pi_{g,\nu}) \arrow[r] \& \operatorname{Autext}(\Pi_{g,\nu}) \arrow[r] \& 1
\end{tikzcd}}
\]

\deleted{\illegible{1}} where the central vertical arrow \deleted{\unsure{Theorem}} is obtained by associating to $g\in\Gamma'_{g,\nu+1}$ the restriction to $\Pi_{g,\nu}\subset\Gamma'_{g,\nu+1}$ of the inner automorphism \unsure{$\operatorname{int}(g)$}.
\mspage{220}

\uline{Theorem \unsure{well known}}\inserted{In the anabelian case${}^*$ $(2g+\nu\geq3)$}
\[
\Gamma_{g,\nu}\xrightarrow{\sim}\operatorname{Autext}_{\mathrm{lac}}(\Pi_{g,\nu}),
\]
Which amounts to the same, \inserted{by (44)}
\[
\Gamma'_{g,\nu+1}\xrightarrow{\sim}\operatorname{Aut}_{\mathrm{lac}}(\Pi_{g,\nu}).
\]
the \unsure{image} of $\Gamma_{g,\nu}\longrightarrow\operatorname{Autext}\Pi_{g,\nu}$ consists of the outer automorphisms which respect the loop structure of $\Pi_{g,\nu}$ of type $g,\nu$.

\margin{${}^*$ \uline{N.B.} In fact this remains true \deleted{in the cases} for the \unsure{pair|case} $(1,0)$ (cf. below) and even in the case $(0,2)$ if we define ad hoc the notion of group with \unsure{loops} $(0,2)$.}

\scanpage{222}
\mspage{111}
(vacuous condition if $\nu=0$, for that matter\ldots).

\uline{Corollary} In the anabelian case $(2g+\nu\geq3)$, the functor $X\longmapsto\pi_1(X)$ from the \inserted{\uline{isotopic}} category of surfaces of type $g,\nu$ to the \uline{outer} groups with loops of type $(g,\nu)$, is an equivalence of categories, as is the functor $(X,s)\longmapsto\pi_1(X,s)$ \inserted{\uline{from the isotopic category}} of \uuline{\unsure{pointed|punctured}} surfaces of type $(g,\nu)$, to the groups with loops of type $(g,\nu)$.

\margin{\uline{N.B.} The arrows in \unsure{the|the} isotopic category are the isotopy classes of \uline{homeomorphisms}.}
\margin{It is even an equivalence at the level of $\infty$-isotopic categories which \unsure{expresses} only the fact that \deleted{\illegible{1}} $A^0_{g,\nu}$ and $A^0_{g,\nu+1}$ are contractible.}

It remains to examine to what extent we can adapt these results to the ``abelian'' case. \deleted{Let us note} \inserted{\uline{Let us recall}} that in this case, we have
\[
(45)\qquad \Gamma'_{g,\nu+1}\xrightarrow{\sim}\Gamma_{g,\nu\deleted{\unsure{+1}}}
\qquad\begin{array}{l}
\text{if }2g+\nu\leq2\\
\text{(``abelian'' cases)}
\end{array}
\]
thus $\Gamma_{g,\nu}$ identifies with the subgroup of $\Gamma_{g,\nu+1}$ formed by the elements which fix $a_{g,\nu}$.

1) Case $g=1$, $\nu=0$ thus $\Gamma_{1,0}\simeq\Gamma'_{1,1}=\Gamma_{1,1}$.

Let us consider the canonical homomorphism
\[
(46)\qquad \Gamma_1=\Gamma_{1,0}\longrightarrow\operatorname{Autext}(\Pi_{1,0})=\operatorname{Aut}(\Pi_{1,0})
\]
\[
(\simeq\operatorname{Gl}(2,\mathbb{Z})\text{ when we have chosen a basis of }\Pi_{1,0}\simeq\mathbb{Z}^2)
\]

which corresponds to the homomorphism
\[
(46\text{ bis})\qquad\Gamma_{1,1}\longrightarrow\operatorname{Aut}\pi_1(U_{1,0},a_{1,0})
\qquad(\simeq\operatorname{Gl}(2,\mathbb{Z})),
\]
deduced \inserted{both of them} by passage to the quotient from the \deleted{\illegible{1}} obvious operations of $A_{1\deleted{\unsure{0}}}=A_{1,0}$\inserted{$=\operatorname{Aut}(X_1)$}, and of
\mspage{221}

\scanpage{223}
$A_{1,1}=\operatorname{Aut}(\unsure{X_{1,1}|U_{1,1}},a_{1,0})$.
It is immediate that this homomorphism is surjective, but less obvious that it is an isomorphism, i.e. that \unsure{every} homeomorphism of $X_1$ which induces the identity \deleted{on} the \deleted{marked point} \inserted{$\Pi_1$ is isotopic} \deleted{$a_{1,0}$ \illegible{1}} to the identity --- this is nevertheless a true (and known) result.

In the present case where \deleted{\illegible{1}} $X_1$ \deleted{\illegible{2}} \inserted{identifies} with the underlying topological space of a topological group $H$ (which is the case in all the abelian cases except that of $g=0$, $\nu=0$ of the $2$-sphere) \deleted{\illegible{1}}, by \unsure{translations|translation} we have a natural homomorphism
\[
(47)\qquad H\longrightarrow A_{g,\nu}\simeq\operatorname{Aut}_{\mathrm{top}}H
\]

allowing us to identify $H$ with a \inserted{topological} subgroup of $A_{g,\nu}$, and we find that the map
\[
(48)\qquad (g,u)\longmapsto gu\qquad
H\times\underbrace{A_{g,\nu+1}}_{\inserted{\begin{array}{c}\text{Stabilizer}\\\text{of }1_H\text{ in}\\ \longleftarrow A_{g,\nu}\simeq\operatorname{Aut}_{\mathrm{top}}(H)\end{array}}}
\longrightarrow A_{g,\nu}
\]
is a homeomorphism. \deleted{\unsure{As}} This \deleted{implies} gives (45) again (since $H$ is connected)
\mspage{222}

\scanpage{224}
\mspage{112}
and makes it considerably more precise, by
\[
(49)\qquad\pi_i(A_{g,\nu+1})\simeq\pi_i(A_{g,\nu})\times\pi_i(H)
\]
\deleted{(\illegible{1} $(g,\nu)$ is \illegible{1} \unsure{trivial}}
Non-spherical abelian cases, i.e. \unsure{any|only} of the three cases $(g=1,\nu=0)$, $(g=0,\nu=1)$, $(g=0,\nu=2)$.

In the ``limiting'' cases $(1,0)$ and $(0,2)$ (when $(g,\nu)$ is abelian, $(g,\nu+1)$ anabelian), since $A^0_{g,\nu}\simeq H\times A^0_{g,\nu+1}$ with $A^0_{g,\nu+1}$ contractible, we find again that $H\longrightarrow A^0_{g,\nu}$ is a homotopism, \deleted{\illegible{2}}

(Let us note that in the cases envisaged\inserted{$(1,0)$, $(0,2)$ and $(0,1)$}, if we choose a complex structure on \unsure{the|the} \unsure{pure} compactification $\widehat X$ of $X=U_{g,\nu}$ and choose $a_{g,\nu}$ as origin, there is a unique $\mathbb{C}$-\uline{algebraic} group structure on $\unsure{\widehat X|X}$ \unsure{admitting} $\unsure{1}$ as unit, and \deleted{\illegible{2}} \inserted{the} \deleted{\illegible{3}} identity component of the group of automorphisms of the algebraic variety $X$ is precisely none other than the group of translations, in the limiting cases $(1,0)$ and $(0,2)$.
\mspage{223}

\scanpage{225}
2) Case $g=0$, $\nu=2$, thus $H\simeq\mathbb{C}^*$,
\[
A_{g,\unsure{\nu|2}}\simeq H\times A_{g,\nu+1}\quad\text{i.e.}\quad A_{g,2}\simeq\mathbb{C}^*\times\uline{A_{g,3}}
\]
thus
\[
(50)\qquad\left\{\begin{array}{l}
\Gamma_{g,2}\simeq\Gamma'_{g,3}\\
A^0_{g,2}\xleftarrow{\simeq}\mathbb{C}^*\quad\text{(homotopy equivalence)}
\end{array}\right.
\]

\deleted{\unsure{As}} $\Pi_{\text{\textcircled{1}},2}\simeq\mathbb{Z}$\deleted{\illegible{1}}, let us consider
\margin{(rather $\Pi_{0,2}$)}
\[
(51)\qquad\Gamma_{g,2}\xrightarrow{\rho}\operatorname{Aut}\left(\underset{\inserted{\begin{array}{c}\Vert\\\pi_1(U_{0,2})\end{array}}}{\Pi_{\text{\textcircled{1}},2}}\right)\deleted{=\pi_1(\unsure{X_{0,2}})}\simeq\{\pm1\}\simeq\mathfrak S_2
\]
identifying with
\[
(52)\qquad\Gamma'_{g,3}\longrightarrow\operatorname{Aut}\left(\underset{\begin{array}{c}\Vert\\\pi_1(U_{0,2},a_{0,2})\end{array}}{\Pi_{\text{\textcircled{1}},2}}\right)\simeq\{\pm\}
\]

This time the homomorphism (51) is \uuline{not} bijective --- it identifies with the canonical homomorphism
\[
\deleted{\unsure{\Gamma_g}}\ (51\text{ bis})\qquad\Gamma_{g,2}\xrightarrow{\rho}\mathfrak S_2
\]
with kernel $\Gamma^!_{g,2}$ --- and $\Gamma^!_{g,\unsure{\nu|2}}$ is not reduced to $1$, for example (if we take $U_{0,2}=\mathbb{C}^*$)

$z\longrightarrow\overline z^{\deleted{\unsure{-1}}}$ defines an element of $\Gamma^!_{g,2}$ which is not \unsure{equal} to $1$. Let us note \deleted{\illegible{1}} \inserted{now} the canonical homomorphism (defined without restriction on $(g,\nu)$)
\mspage{224}

\scanpage{226}
\mspage{113}
\[
(53)\qquad\Gamma_{g,\nu}\xrightarrow{\chi}\{\pm1\}
\]
\margin{This homomorphism is always surjective. We shall denote \unsure{its kernel} by $\Gamma^+_{g,\nu}$ (and no \inserted{longer} $\Gamma^0_{g,\nu}$!).}

via the action of $\Gamma_{g,\nu}$ on the orientation module $T_g=T$ of $\unsure{U_{g,\nu}}$ \deleted{\illegible{1}} (which is defined, for $(g,\nu)\neq(0,0),(0,1)$, \inserted{$(0,2)$}, in terms of the loop structure of $\Pi_{g,\nu}$). In the present case, putting together $\rho$ and $\unsure{\rho|\chi}$, we find a homomorphism

\[
(54)\qquad g\longmapsto(\rho(g),\chi(g)):\quad\Gamma_{g,2}\longrightarrow\mathfrak S_2\times\{\pm1\}
\simeq\deleted{\unsure{\mathbb{Z}_2}}(\pm1)\times(\pm1)
\]
which is obviously surjective (if $g$ corresponds to $z\longmapsto z^{-1}$, its image is \unsure{$(-1,-1)$|$(-1,1)$}, if it \unsure{corresponds to} $z\longmapsto\overline z$, its image is $(1,-1)$).

\margin{\uline{N.B.} This suggests that for an isotopic description of the category of surfaces \unsure{$X$} of \inserted{type} $(0,2)$, we must use the pair of $\pi_1(X)$ $(\simeq\mathbb{Z})$ and its orientation module $T$; \unsure{likewise|ditto} below concerning \unsure{the cases of types|the case of type} $(0,1)$, but then $\pi=0$ and $T$ suffices.}

\unsure{I} say that it is \uline{bijective}, \deleted{\illegible{3}}
\inserted{$\displaystyle\Gamma_{g,2}\xrightarrow[\rho,\chi]{\sim}\mathfrak S_2\times\pm1$}
which amounts to the same, that the restriction of $\rho$ (51 bis) to the kernel \inserted{$\Gamma^+_{0,2}$} of $\chi:\Gamma_{g,2}\longrightarrow\{\pm1\}$ is an isomorphism, or what amounts to the same

\uline{Theorem} (well known!)
\[
(55)\qquad\Gamma^{!+}_{0,2}=1,
\]
or again:
$\rho^+_{0,2}:\Gamma^+_{0,2}\longrightarrow\mathfrak S_2$ is an \uline{isomorphism}, or $\chi^!_{0,2}:\Gamma^!_{0,2}\xrightarrow{\unsure{\sim}}\pm1$.

[\uline{N.B.} If we want a statement common to the \deleted{cases} two ``limiting abelian'' cases $(1,0)$ and $(0,2)$, we shall say that in these cases, $\Gamma^+_{g,\nu}\longrightarrow\operatorname{Aut}(\Pi)$
\mspage{225}

\scanpage{227}
is \uline{injective}, and has as image \deleted{\unsure{SL}} the group of automorphisms of the free $\mathbb{Z}$-module $\Pi$ (of rank $2$ or $1$) which are of \uline{determinant $1$}.]

This is equivalent, modulo the isomorphisms $\Gamma'_{g,3}\simeq\Gamma_{g,2}$, to the \deleted{\illegible{1}}

\uline{Corollary}
\[
(56)\qquad\left\{\begin{array}{l}
\Gamma^+_{\unsure{g|0},3}\xrightarrow{\rho^+_{0,3}}\mathfrak S_3\quad\text{is an }\uline{\text{isomorphism}},\text{ and}\\
\Gamma_{0,3}\xrightarrow{(\rho_{0,3},\chi_{0,3})}\mathfrak S_3\times(\pm1)\quad\text{also}
\end{array}\right.
\]

\deleted{As for the whole of $\Gamma_{0,3}$, it is an extension of \deleted{\illegible{1}} $\{\pm1\}$ by $\Gamma^+_{0,3}$, which of course can be split (\uline{N.B.} This is so for $\Gamma_{g,\nu}\xrightarrow{\chi}\{\pm1\}$ whatever $(g,\nu)$), and \unsure{to be determined|is determined} easily as the \unsure{product} group.}

3) Case $g=0$, $\nu=1$ thus $H\simeq\mathbb{C}$, thus
\deleted{$A_{g,1}$}
\[
(57)\qquad A_{g,1}\overset{\text{homeomorphism}}{\simeq}\mathbb{C}\times A_{g,2}\simeq\underbrace{\mathbb{C}\times\mathbb{C}^*}_{\operatorname{Aff}(1,\mathbb{C})}\times A_{g,3}
\]
\inserted{(which is \unsure{simply transitive} on pairs of distinct elements)}
\deleted{\unsure{whence}}

\margin{N.B. Since $A^0_{g,3}$ is contractible, this indeed \unsure{shows again|shows} that the inclusion of $G=\operatorname{Aff}(1,\mathbb{C})=\operatorname{Aut}_{\mathbb{C}}(X)=\operatorname{Aut}_{\mathbb{C}}(X)^0$ into $A_{g,1}$ is a homotopism.}

taking on $\unsure{U_{g,1}}$ a complex algebraic group structure $\simeq\mathbb{C}$.
\mspage{226}

\scanpage{228}
\mspage{114}
We have \inserted{here} $\Gamma_{g,1}\simeq\Gamma'_{g,2}=\Gamma^!_{g,\unsure{1|2}}\xrightarrow[\sim]{\chi}(\pm1)$

i.e.

\uline{Corollary} We have an isomorphism via the sign
\[
(58)\qquad\Gamma_{\unsure{0|g},1}\xrightarrow[\sim]{\chi_{\unsure{0|g},1}}\{\pm1\}
\]

4) Case $g=0$, $\nu=0$. Here there is no $H$ i.e. no topological group structure on $X\simeq\deleted{\unsure{S^2}}X_0=U_{0,0}\simeq S^2$, but taking a complex structure on $X$ (whence $X\overset{\inserted{\mathbb{C}\text{-isomorphism}}}{\simeq}\mathbb{P}^1_{\mathbb{C}}$) we find a group $G$ of \deleted{\unsure{homeomorphisms}} \inserted{$\mathbb{C}$-automorphisms} $G\simeq\operatorname{GP}(1,\mathbb{C})$, and
\[
(59)\qquad G\hookrightarrow A_0=A_{00},
\]

\unsure{let us prove} again that it is an \deleted{equivalence} of homotopism. Since $G$ is simply transitive \unsure{on|\deleted{on}} the \deleted{systems} \inserted{triples} of $3$ distinct points of $X_0$, we again find a homeomorphism
\[
(60)\qquad G\times A^!_{0,3}\longrightarrow A_0\qquad(g,u)\longmapsto gu
\]

whence
\[
(61)\qquad\Gamma_0\simeq\Gamma^!_{03}\xrightarrow[\sim]{\chi_{0,3}}\{\pm1\}
\]
\mspage{227}

\scanpage{229}
and taking into account that the identity component $A^0_{0,3}$ of $A^!_{0,3}$ is contractible, we indeed find \inserted{again} that (59) is an \deleted{equivalence} of homotopism.
\begin{center}\rule{1.2cm}{0.25pt}\\[-0.8em]\rule{1.2cm}{0.25pt}\end{center}

\uline{Common conclusion} \inserted{to all cases.} It is appropriate to include in the notion of ``group with loops'' also the \inserted{four} abelian cases, and we do so as follows:

1) Type $(1,0)$ --- we do not have to supplement the general definition, \inserted{which here amounts to saying that} $\Pi$ is an abelian group, a free $\mathbb{Z}$-module of rank $\unsure{1|2}$. Its orientation module $T$ can indeed be defined as $\unsure{H_2}(\Pi,\mathbb{Z})$, or as $\Lambda^2_{\mathbb{Z}}\Pi$. The \deleted{\unsure{action}} \inserted{sign} of an automorphism of $\Pi$ is \unsure{given|given} by its action on $T$, it is also its determinant.

2) Type $(0,2)$, $\Pi\overset{\inserted{\text{not canonical}}}{\simeq}\mathbb{Z}$. Here we must give ourselves, \uline{in addition to} $\Pi$, an invertible $\mathbb{Z}$-module $T$, and the loop structure is defined by one of the two isomorphisms $T\simeq\mathbb{Z}$. An automorphism of the loop structure is therefore defined by \deleted{a} \inserted{(any} pair \deleted{\illegible{1}} $u_{\Pi},u_T$ of an automorphism of $\Pi$ and one of $T$.
\mspage{228}

\scanpage{230}
\mspage{115}
3 and 4) Types $(0,1)$ and $(0,0)$, which are those where $\Pi=\pi_1(U_{g,\nu})$ is zero. By definition, a structure of group with loops of type $(g,\nu)$ is defined here by the datum of an ``orientation module'' and nothing more, which is an invertible $\mathbb{Z}$-module $T$, the automorphisms are those of $T$.
\margin{and we must give \unsure{it|ourselves} in addition the type, i.e. (with $g=0$ understood) the $\nu\in\{0,1\}$.}

With these definitions, and for fixed $g,\nu$, we have the

\uline{Theorem} \inserted{Let $(g,\nu)\in\mathbb{N}\times\mathbb{N}$.} The natural functor from the \inserted{\uline{isotopic}} category of surfaces of type $(g,\nu)$ to \inserted{the category} \unsure{of} outer groups with loops of type $(g,\nu)$, is an equivalence \unsure{of|of the} categories. (N.B. as arrows we take the classes of isomorphisms, on one side as on the other).

But this theorem is of course not fully satisfactory in \unsure{the abelian case|the abelian cases}, for example the datum of an object of the isotopic category (\unsure{made explicit|made explicit} by \unsure{its|an} outer $\pi_1$ with loops), in the case of an action of a group, does not even allow us to recover
\mspage{229}